\documentclass[11pt,a4paper]{article}
\usepackage[T1]{fontenc}
\usepackage[margin=1in]{geometry}
\usepackage{amsmath,amssymb,amsthm,mathtools}
\usepackage{booktabs}
\usepackage{xcolor}
\usepackage[colorlinks=true,linkcolor=blue!60!black,citecolor=blue!60!black,
            urlcolor=blue!60!black]{hyperref}
\usepackage{microtype}
\usepackage{enumitem}
\usepackage[most]{tcolorbox}

\newcommand{\lean}[1]{\texttt{\small #1}}
\newtcolorbox{keybox}[1][]{colback=blue!3,colframe=blue!45!black,
  boxrule=0.5pt,left=4pt,right=4pt,top=3pt,bottom=3pt,#1}
\newtcolorbox{failbox}[1][]{colback=red!3,colframe=red!45!black,
  boxrule=0.5pt,left=4pt,right=4pt,top=3pt,bottom=3pt,#1}
\newtheorem{theorem}{Theorem}

\title{\textbf{Duality is a structure; the sector is the cause}\\
\large T-duality of the compact boson stated on its sectors and checked by two kernels,\\
the measured invariant it predicts, and a criterion for when a self-dual point is physical}

\author{Xavier Callens\\
\small SocrateAI-Scientific-QuantumFluids\\
\small \url{https://github.com/xaviercallens/SocrateAI-Scientific-QuantumFluids}}
\date{September 2026}

\begin{document}
\maketitle

\begin{abstract}
\noindent
An involution $R\leftrightarrow\alpha'/R$ on a length scale was, in earlier work of this programme, offered a
physical home in T-duality, in K3 mirror symmetry, and in the modular symmetry of the quantum Hall system. Three
rounds of pre-registered numerics and twenty-eight machine-checked theorems later, we can say where such a
duality lives and what it does. We state T-duality on the sectors $(n,w)\in\mathbb Z^2$ of the compact boson and
prove, in Lean~4 against Mathlib with two independent kernels, that (i) it is a reindexing of the sector lattice
that leaves the sector sum invariant with no convergence hypothesis, (ii) the product of the electric and magnetic
scaling dimensions is radius-independent, $\Delta_e\Delta_m=\tfrac14$, which in the superfluid dictionary is
$\eta\,n_s\lambda_T^2=1$ -- the relation this programme measured to $12$--$27\%$ on eleven energies in two boxes --
and (iii) the self-dual radius $\sqrt2$ and the radius $2\sqrt2$ at which the vortex becomes marginal are
different points: the Berezinskii--Kosterlitz--Thouless transition is the threshold of one sector, not the fixed
point of the exchange. From these three facts and the literature on Ising, string-gas and superconductor--insulator
dualities we extract a criterion: a self-dual point pins a transition only when the dual sectors are simultaneously
critical on a discrete set of fixed points; on a line of fixed points it is a point like any other; when both
sectors are populated it can be a preferred equilibrium. In every case what acts is a sector -- a conserved label
that crosses a threshold, proliferates, annihilates, or is imposed -- and the duality relabels it. The duality is
exact, verified, and not a cause. No physics here is new; the mathematics is elementary; what is new is that the
distinction is stated so that a kernel can check it and an experiment can test it.
\end{abstract}

\section{Three questions one involution raised}

The dual-length involution $k\mapsto k_s^2/k$ of~\cite{Callens2026Lean} has the form of T-duality. It was pointed
out that the same involution is the Fricke element of $\Gamma_0(N)^+$, the mirror involution of degree-$2N$
lattice-polarised K3 surfaces~\cite{Dolgachev1996}, and that $\Gamma_0(2)$ acts on the quantum Hall conductivity
plane~\cite{LutkenRoss1992}. Each identification was taken to its gate~\cite{Callens2026Wasserstein,
Callens2026Causal}: the Fricke matrix normalises Mathlib's $\Gamma_0(n)$ and its axis action is the dual-length map,
but no $\Gamma_0(n)$ acts on a wavenumber (\lean{Fricke.lean}); the Fricke element fixes the Hall critical point's
orbit but sends every plateau off the odd-denominator class (\lean{QHFricke.lean}); and the physical BKT point is not
the self-dual radius. The residue of these three verdicts is one question: if a duality is exact and verified, and
yet is not where the physics happens, what is it, and what is?

\section{T-duality on the sector lattice, machine-checked}

Let the compact boson have radius $R$ in the normalisation where the self-dual radius is $\sqrt2$. The primary of
momentum $n$ and winding $w$ has $h=\tfrac12(n/R+wR/2)^2$, $\bar h=\tfrac12(n/R-wR/2)^2$, hence
\[
\Delta(n,w;R)=n^2/R^2+w^2R^2/4,\qquad s=h-\bar h=nw .
\]
The electric (spin-wave, $e^{i\theta}$) operator is $(1,0)$; the magnetic (vortex) operator is $(0,1)$. The
following are theorems of \lean{CompactBoson.lean} (twelve statements, axiom footprint $\{\mathrm{propext},
\mathrm{Classical.choice},\mathrm{Quot.sound}\}$, accepted by the Lean kernel and by \lean{nanoda} through
Comparator; three negative controls -- product $\tfrac12$, BKT at $\sqrt2$, duality with $R\mapsto1/R$ -- fail).

\begin{theorem}[duality as reindexing]
$\Delta(n,w;R)=\Delta(w,n;2/R)$ and $s(n,w)=s(w,n)$. Consequently, for every weight $F$,
$\sum_{(n,w)\in\mathbb Z^2}F(\Delta(n,w;R),s)=\sum_{(n,w)\in\mathbb Z^2}F(\Delta(n,w;2/R),s)$.
\end{theorem}
The second statement is \lean{Equiv.tsum\_eq} applied to the swap of $\mathbb Z^2$: reindexing a sum by a
bijection is an identity of \lean{tsum} whether or not the series converges. This is what the duality is -- a
bijection of the sector lattice that preserves the spectrum. Savit's review~\cite{Savit1980} gives the general
form: a duality is a Fourier transform on the group of sectors, electric charges in $G$ against magnetic charges
in $\hat G$, with the Lagrangian sum over one family turned into the Hamiltonian sum over the dual family by Poisson
summation.

\begin{theorem}[the invariant]
$\Delta(1,0;R)\,\Delta(0,1;R)=\tfrac14$ for every $R\neq0$.
\end{theorem}
In the XY / two-dimensional-superfluid dictionary $\Delta(1,0)=\eta/2$ and $\Delta(0,1)=n_s\lambda_T^2/2$, so the
invariant reads $\eta\,n_s\lambda_T^2=1$. That is the relation measured in~\cite{Callens2026Causal}: $1.047$ and
$1.072$ on two equilibrated energies in round~2, $1.12$--$1.22$ on the six energies of the heating ladder,
$1.14$--$1.27$ on the five energies of the $L=32$ box -- with the finite-size excess predicted by
Hasenbusch's logarithmic corrections~\cite{Hasenbusch2005}. \emph{The numerical verification of the duality in
this programme is the measurement of this product.}

\begin{theorem}[self-dual is not marginal]
For $R>0$: $\Delta(1,0)=\Delta(0,1)$ iff $R=\sqrt2$, where both equal $\tfrac12$; $\Delta(0,1)=2$ iff
$R=2\sqrt2$, where $\Delta(1,0)=\tfrac18$, i.e.\ $\eta=\tfrac14$; and $2\sqrt2\neq\sqrt2$.
\end{theorem}
The transition is where the vortex operator becomes marginal -- Nelson and Kosterlitz's universal jump
$n_s\lambda_T^2=4$~\cite{NelsonKosterlitz1977}, located in our own runs at $T_{\rm BKT}=0.821$ with
$\eta=0.31$ there -- and that is a condition on one sector. The exchange symmetry has its fixed point elsewhere,
and nothing happens there.

\section{When is a self-dual point physical?}

The compact boson is the case with a \emph{line} of fixed points, one for each $R$. Three other cases complete the
criterion.

\paragraph{Ising: yes.} Kramers--Wannier self-duality~\cite{KramersWannier1941} fixes $T_c$ at $\sinh2K=1$ because
the order and disorder operators~\cite{FradkinSusskind1978} exchange under the duality and $\mathbb Z_2$ has a single
symmetric point: the dual sectors are simultaneously critical and the fixed-point structure is discrete.

\paragraph{String gas cosmology: a preferred equilibrium.} In the Brandenberger--Vafa
scenario~\cite{BrandenbergerVafa1989,Brandenberger2008} ``the symmetry we make use of is T-duality, and the new
degrees of freedom are the string oscillatory modes and the string winding modes''. Momentum modes, with energy
$n/R$, resist contraction; winding modes, with energy $wR$, resist expansion; when both are populated the radion is
stabilised at the self-dual radius as the minimum of an effective potential. The self-dual point is physical there
because both dual sectors are present and compete. And the exit from the Hagedorn phase -- the event that makes
three dimensions large -- is the \emph{annihilation of winding modes}: a sector event, with T-duality as the
symmetry of the spectrum in which it is described.

\paragraph{Superconductor--insulator transition: a constraint.} Boson--vortex duality~\cite{DasguptaHalperin1981,
FisherLee1989} predicts, if the transition is self-dual, a universal resistance $h/4e^2$~\cite{Fisher1990,
ChaFisherGirvin1991}; measured values scatter around it. The cause of the transition is the proliferation of
vortices on one side and of charges on the other; self-duality constrains the critical point only if the fixed
point happens to be self-dual, which is a property of the fixed point, not of the duality.

\paragraph{Quantum Hall: organisation, not cause.} $\Gamma_0(2)$ organises the plateau diagram and the position
of the critical points~\cite{LutkenRoss1992}; the cause of the quantisation is Laughlin's gauge
argument~\cite{Laughlin1981} on the Chern sector; and the Fricke element, which fixes the critical orbit, is not a
symmetry of the plateaux (\lean{QHFricke.lean}).

\begin{keybox}
\textbf{Criterion.} A self-dual point pins a transition iff the two dual sectors are simultaneously critical and
the fixed-point set is discrete. On a line of fixed points the self-dual point is a point like any other and the
transition sits where one sector crosses its own threshold. When both dual sectors are populated and compete
energetically, the self-dual point is a preferred equilibrium. In each case the duality relabels the sectors and
the sector acts.
\end{keybox}

\section{What acts: the sector}

The same triad recurs in every case above: a group of sectors ($\mathbb Z$ for $U(1)$ windings, $\mathbb Z_2$ for
Ising, Chern numbers, string windings); a duality that is a Fourier or reindexing map between $G$-sectors and
$\hat G$-sectors preserving the spectrum; and physics that happens when a sector crosses a threshold, proliferates,
annihilates, or is imposed. The causal role of the sector is what the rest of this programme
established~\cite{Callens2026Causal}: it is conserved by every continuous evolution without a phase slip
(\lean{TopologicalProtection}), additive in scale so that a bound pair is invisible from outside
(\lean{ScaleResolvedWinding}), read correctly by the sampled loop sum as the continuum degree
(\lean{ContinuumWinding}), it carries its own temperature (\lean{SectorTemperature}), and in an energy-matched
intervention it is a difference-maker: the same energy delivered as topology lowered the condensate by
$0.49$--$0.73$ where phonons lowered it by at most $0.035$, with the injected sector reading five to eight times
the bath's temperature on our own calibrated thermometer. That is why one involution appears in helium, Hall bars,
K3 moduli and string gases, and why it explains none of them.

\section{What this paper does not claim}
\begin{failbox}
\begin{itemize}[leftmargin=*]
\item Novelty in the mathematics: the three theorems are elementary; their value is that the distinction between
  the duality and the sector is stated so that a kernel checks it.
\item Novelty in the physics: Kramers--Wannier, Nelson--Kosterlitz, Brandenberger--Vafa, Fisher and Lütken--Ross
  are cited, not extended.
\item Any identification of a superfluid with a K3 surface or a string background. The K3 mirror involution is the
  same map on a different space; the criterion above is what the shared map does and does not license.
\item That every duality in physics fits the triad: the cases examined do; the general statement is a reading.
\end{itemize}
\end{failbox}

\bibliographystyle{unsrt}
\bibliography{refs_duality_sector}
\end{document}
